\section{Conclusiones}

Se implementó el juego de reflejos completo sobre el Pochoco SoC: temporización de tres
segundos, elección pseudoaleatoria del LED objetivo, medición del tiempo de respuesta
en décimas de segundo, manejo de botones incorrectos y promedio de diez rondas. Todo
el comportamiento se verificó en simulación sobre el RTL del SoC y sobre el circuito
sintetizado, el diseño sintetiza para la iCE40 HX1K con margen de frecuencia, y el juego
funciona en la Go Board.

El único cambio de hardware necesario fue un contador de ciclos de 32 bits. Con él, toda
medición de tiempo se reduce a una resta sin signo, que además tolera que el contador
dé la vuelta. La alternativa, bucles de espera calibrados, habría acoplado la duración
de cada fase del juego al número exacto de instrucciones del programa.

Las restricciones del procesador obligaron a expresar operaciones elementales con las
instrucciones disponibles: multiplicar por dos es sumar un número consigo mismo, leer
el bit 31 es preguntar si el número es negativo y dividir es restar repetidamente. El
mismo criterio permitió implementar un LFSR sin desplazamientos. El modelo del
generador mostró que el número de pasos entre elecciones importa: con pocos pasos el
primer LED de cada partida era siempre el mismo, aunque la semilla viniera del
contador.

La prueba en la placa mostró el límite de la simulación. El juego pasaba todas las
simulaciones, incluida la del circuito sintetizado, y aun así no arrancaba en la FPGA,
porque el reset de encendido original soltaba al procesador antes de que la RAM
estuviera disponible. Los modelos de simulación no representan ese retraso. La causa
se aisló grabando programas de prueba cada vez más cercanos al juego y descartando una
hipótesis por vez, y la corrección fue una línea de Verilog.

En el assembler, el problema de fondo son las referencias hacia adelante: una etiqueta
se usa antes de saber su dirección. Se resolvió con pasadas separadas para las
direcciones y para la codificación. La decisión que más confianza dio sobre su
corrección fue probarlo contra referencias externas: la salida del assembler de GNU,
las fórmulas del decodificador del procesador y el procesador mismo, en simulación.
